% !TeX root = ../main.tex

% TODO(author): 中文摘要為英文摘要之翻譯初稿，請核對術語與行文後定稿。
\begin{abstract}

大型語言模型（LLM）已能生成語法正確的程式碼，然而在真實硬體專案中，暫存器傳輸級
（RTL）設計工作多以整合與矽智財（IP）重用為主，而非從零撰寫演算法模組。本論文以
RealBench 基準測試中源自三個真實開源 IP 專案（AES 加密核心、SD 卡控制器、E203
RISC-V 處理器核心）的 60 個模組級任務為評測對象，研究如何以系統工程手段——規劃、
接地（grounding）、驗證、修復與路由——提升固定開源權重模型的 RTL 生成能力。

本論文提出兩項主要貢獻。其一為兩階段代理式 IP 重用流水線：規劃代理透過
IP 重用工具（目錄檢索、候選 IP 檢視與評分、介面相容性檢查）針對每個任務
的真實 IP 目錄進行規劃，並以「建構式接地」（目錄注入、封閉詞彙表對齊、
完整性檢查）產生可稽核的結構化重用計畫，再由生成器在忠實鏡射評分環境的
Verilator 驗證與修復迴圈下產生 Verilog。在溫度為零的對照實驗中，流水線通過率達 25.0\%（語法 61.7\%），
為同一模型直接提示（10.0\%／26.7\%）的 2.5 倍，且增益完全集中於重用密集的
E203 家族。其二為兩層串接路由器，依任務性質在完整流水線與低成本直接生成之間
逐任務選擇流程：在單變量消融中，解題數由 21/60 提升至 34/60，總運算時間反而由
213 小時降至 130 小時，每解一題的 token 效率為暴力混合策略的約 4 倍。

本論文同時以對照消融精確量化各輔助機制：語意修復快取單調提升語法率
（59.4\%→66.9\%），並使通過樣本更集中於可解任務；兩種修復階段皆呈
「速修或永不修」分布，精簡的修復預算（語法約 6 次、功能不超過 4 次）
即可保留幾乎全部的修復成功案例，同時大幅節省運算；以黃金模型稽核法
檢驗持續零分任務，發現並向上游回報一項測試平台競爭條件（已獲基準採納
修復），亦確認其餘零分為真實難度。整體結果顯示：下一步的突破點不在
語法、規劃或模型規模，而在功能正確性——需要比失配計數更豐富的回饋
訊號。

\end{abstract}

\begin{abstract*}

Large language models generate syntactically plausible code, but real-world
RTL engineering is dominated by integration and IP reuse rather than
freestanding algorithm writing. This thesis studies how far system
engineering --- planning, grounding, verification, repair, and routing
around fixed open-weight models --- can push LLM-based Verilog generation
on RealBench, a benchmark of 60 module-level tasks drawn from three real
open-source IP projects (an AES cipher core, an SD-card controller, and
the E203 RISC-V processor core) and scored by a strict warnings-as-fatal
compile-and-simulate harness.

The thesis makes two primary contributions. First, a two-stage agentic
IP-reuse pipeline: a planning agent works each task through a suite of
IP-reuse tools --- catalog search, candidate inspection and scoring, and
interface-compatibility checking --- and produces a structured, auditable
reuse plan grounded by construction against a per-task catalog of the
modules the environment actually provides (catalog injection,
closed-vocabulary alignment, and a completeness gate), followed by a
generator operating under a Verilator verify-and-repair loop that
mirrors the scoring environment exactly. At temperature zero the pipeline passes
25.0\% of tasks (61.7\% syntax) against the same model's direct-prompt
10.0\% (26.7\%), a $2.5\times$ improvement concentrated entirely in the
reuse-rich processor-core family. Second, a two-tier cascade router that
matches task character to generation flow, sending integration-shaped
tasks to the pipeline and self-contained algorithmic tasks to a cheap
direct flow: in its controlled ablation it raises solved tasks from 21 to
34 of 60 while reducing total compute from 213 to 130 hours, reaching
four times the token efficiency per solved task of brute-forcing both
flows.

Controlled ablations quantify the supporting machinery with equal care:
a semantic repair cache lifts syntax monotonically (59.4\%$\to$66.9\%)
while concentrating passing samples inside solvable tasks; both repair
phases converge within their first attempts, so compact budgets
($\approx$6 syntax, $\leq$4 functional) retain essentially all rescues
at a fraction of the compute; and a golden-model audit of persistently
failing tasks uncovered a testbench race condition --- since adopted and
fixed upstream by the benchmark --- while confirming the remaining zeros
as genuine difficulty. The clearest opportunity at the end is neither
syntax, nor planning, nor model scale, but functional correctness, which
will reward a richer feedback signal than a mismatch count.

\end{abstract*}
